\frontmatterpage
\unchapter{Abstract}

In many countries of the former Soviet Union there are small offices that exchange one currency for another and send money abroad. A family in Uzbekistan, for example, may walk into such an office to receive money sent by a relative working in Russia, and the office, in turn, settles up with a partner office in Moscow once a week. These offices are often family businesses with branches in several cities and countries. They handle a lot of money every day, but they usually keep their records in a paper notebook or in an Excel spreadsheet. As a result, mistakes happen often, and a single missed line can lead to a loss of several thousand US dollars that takes weeks to find.

The reason mistakes happen is that the office really has to keep three kinds of accounts at the same time. The first is the money in its own cash drawer and on its own bank cards. The second is the running debt with every partner office abroad, and this debt has to be tracked in several currencies at once, because the same partner can owe the office both Russian roubles and Turkish liras at the same time. The third is the money belonging to regular customers, who often keep a small balance with the office in different currencies. Standard accounting programs mix all three of these into one number in one currency, and that hides exactly the things the office needs to see.

This Bachelor's Thesis describes a programme called \emph{Ethnocount}, which the author designed and built to solve this real problem. The programme works on phones (Android, iPhone) and on computers (Windows, macOS, Linux) from the same code, so a branch can use the same tool whether the operator is sitting at a desk or running between clients with a phone in hand. The programme keeps the three kinds of accounts separate, the way they really are; it adds money up correctly even when the same partner has debts in three different currencies; and it makes sure that a transfer cannot be saved twice by accident, even if the operator taps the button twice on a bad internet connection. There are also rules about who can do what: the founder of the network sees everything, a director sees all branches, and a regular employee sees only his or her own branch and has to ask the director's permission before changing past records.

The author built the system, tested it carefully with many automatic tests, measured how fast it works on a normal Windows computer and on a middle-class Android phone, and tried it for two weeks together with a real exchange office. During the trial the office kept both its old Excel spreadsheet and the new programme up to date, and the two were compared at the end of every week. The new programme did not introduce any mistakes of its own. On the contrary, it found three old mistakes in the spreadsheet that the operator had not noticed before.

The work still has limits, which are stated honestly. The programme does not yet download exchange rates from the internet on its own; the operator has to type them in. It does not let the operator save a transfer when there is no internet --- in such a case the action is refused, because for money it is safer to refuse than to risk a wrong order later. And it was tested in one office over two weeks rather than in many offices over many months. These limits point to the natural next steps for the project, which are listed in the closing chapter.

\thesisscopeEN

\textbf{Keywords:} cross-platform mobile development, Flutter, BLoC, Supabase, PostgreSQL, Row-Level Security, stored procedures, double-entry accounting, money transfers, foreign exchange, multi-branch operations, multi-tenant architecture, idempotency.
